\medskip

\emph{Pour chacune des affirmations suivantes, indiquer si elle est vraie ou fausse.\\
Chaque réponse doit être justifiée. Une réponse non justifiée ne rapporte aucun point.\\ Dans cet exercice, les questions sont indépendantes les unes des autres.}

\medskip

\begin{enumerate}
\item On considère l'équation différentielle :

\[(E )\quad y' = \dfrac 12y + 4.\]

\smallskip

\textbf{Affirmation 1 :} Les solutions de $(E)$ sont les fonctions $f$ définies sur $\R$ par :

\[f(x) = k\e^{\frac 12 x} - 8, \quad \text{avec}\: \: k \in \R.\]

\item Dans une classe de terminale, il y a 18 filles et 14 garçons.

On constitue une équipe de volley-ball en choisissant au hasard 3 filles et 3 garçons.

\smallskip

\textbf{Affirmation 2 :} Il y a \np{297024} possibilités pour former une telle équipe.
\item Soit $(v_n)$ la suite définie pour tout entier naturel $n$ par :

\[v_n = \dfrac{n}{2 + \cos (n)}.\]

\textbf{Affirmation 3 :} La suite $(v_n)$ diverge vers $+\infty$.
\item Dans l'espace rapporté à un repère orthonormé \Oijk, on considère les points A(1~;~1~;~2),\: B$(5~;~- 1~;~8)$ et C(2~;~1~;~3).

\textbf{Affirmation 4 :} $\vect{\text{AB}} \cdot\vect{\text{AC}} = 10$ et une mesure de l'angle $\widehat{\text{BAC}}$ est $30\degres$.
\item On considère une fonction $h$ définie sur $]0~;~+ \infty]$ dont la dérivée seconde est
définie sur $]0~;~+ \infty]$ par :

\[h''(x) = x \ln x - 3x.\]

\textbf{Affirmation 5 :}  La fonction $h$ est convexe sur $\left[\e^3~;~+\infty\right[$.
\end{enumerate}
